\section{There is no God but God}

In Introduction to the Study of Hindu Doctrines, \textbf{Rene Guenon} includes a chapter on the relationship between theology and metaphysics. He writes:

\begin{quotex}
The theological point of view is but a particularization of the metaphysical point of view … it is an application of it to contingent conditions, the mode of adaptation being determined by the nature of the conditions to which it must respond … From this it follows that every theological truth, by means of a transposition dissociating it form its specific form, may be conceived in terms of the metaphysical truth corresponding ot it, of which it is but a kind of translation. 

\end{quotex}
Since everything has a reason for its existence, so does theology. In particular, man as he is a contingent being, will respond and understand theological language, while finding metaphysical language abstract and unapproachable. The other more important factor, de-emphasized or ignored by Guenon, is that metaphysical doctrine itself is secondary and derivative. Specifically, in a book dedicated to the six orthodox schools of Hinduism, which claim to be based on the Vedas, there is no mention of the contents of the Vedas themselves.

The authority of the Vedas derives from their claim to be a revelation from a higher source to rishis or seers in a state of higher consciousness. The Vedas consist of poems, prayers, descriptions of sacrificial rites and rituals, incantations, laws and so on, all things that Guenon might dismiss as ``sentimental" when they are actually foundational.

As an example we can consider the metaphysical doctrine ``\textbf{Being is}" and its theological equivalent, ``\textbf{God exists}". This equivalence was certainly known in the Middle Ages as we see from Thomas Aquinas\footnote{\url{https://www.gornahoor.net/?p=1299}}. Since what is Not Being, isn't, the clear corollary is that there is only one God. Hence, a commitment to monotheism\footnote{\url{https://www.gornahoor.net/?p=564}}, provided it is properly formulated and understood, is necessary to tradition.

Guenon explicates the metaphysics of Being and nonbeing most fully in \textit{The Multiple States of the Being}. Anyone with a logical mind and sufficient powers of concentration can follow his presentation; a fortiori, there is no requirement to be in a state of higher consciousness. For Guenon, Being is not the Infinite\footnote{\url{https://gornahoor.net/?p=3261}}, which contains all possibilities. Being consists of all possibilities that are manifested. Nonbeing, then, consists of non-manifested possibilities as well as possibilities of non-manifestation.

Now Guenon is somewhat inconsistent in his various works. On the one hand in \textit{Hindu Doctrines}, he makes the rather racist claim that ``\emph{Westerners, including even those who were true metaphysicians up to a certain point, have never known metaphysic in its entirety}." Yet, he also claims that the Middle Ages knew Tradition and had true initiates\footnote{\url{https://gornahoor.net/?p=3281}}. In particular, he writes that Western Neo-Platonism did indeed understand the Infinite. As Gornahoor has pointed out\footnote{\url{https://gornahoor.net/?p=2298}}, the intellect of the Middle Ages was formed by the Neo-Platonists \textbf{Plotinus}, \textbf{Augustine}, and \textbf{Boethius}.

If God is the principle of Being, then how can God also be Infinite? Augustine understood the issue at stake. For Augustine, the possibilities in the Infinite are precisely what Plato means by ideas. Then, Augustine places these ideas in the mind of God\footnote{\url{http://readingthesumma.blogspot.com/2010/05/question-15-ideas-in-mind-of-god.html}}. Plato, and even Guenon, assume the ideas subsist in a domain of their own. This solution poses a dual conundrum. First of all, their metaphysical status is unclear. For Guenon, unmanifested possibilities do not ``exist", otherwise they would be manifested possibilities. Hence, there is no way to ``know" them, although he does claim to know them. Augustine recognizes this problem and solves it by putting the Ideas in the mind of God, so that they appear in consciousness but not in the world of manifestation. Hence, for Augustine God is beyond Being and is identical with the Infinite.

The other conundrum is more subtle and cannot be resolved by doctrine alone. One of the objections to Plato's philosophy is the lack of an adequate explanation as to how the ideas become real or manifested. This is a question we discussed in our interpretation of Evola's \textit{The Individual and the Becoming of the World}\footnote{\url{https://www.amazon.com/dp/0999086049}}. This requires the \textbf{Will}, a conclusion that Augustine also reached; the will is not amenable or reducible to any verbal or intellectual theory.

A final point to be made is the status of the possibilities of non-manifestation. Once, again, we see a complete understanding in the West of these, under the concept of \textbf{privation}\footnote{\url{https://gornahoor.net/index.php?s=privation}}. Rather than the trivial examples given by Guenon, Augustine relates this to the concept of ``evil", which, he claims, cannot have a real existence, but is really a privation of the Good. As for Guenon's examples, the ideas of the Void\footnote{\url{https://gornahoor.net/?p=17095}} and of the Silence\footnote{\url{https://gornahoor.net/?p=16873}} are certainly known in the West.

Next we will discuss how we can know Being or God.



\flrightit{Posted on 2011-11-06 by Cologero }

\begin{center}* * *\end{center}

\begin{footnotesize}\begin{sffamily}



\texttt{Enda Miller on 2011-11-07 at 15:22 said: }

What do you mean by the term ``manifested"? More precisely, what does it mean to say a ``manifested or unmanifested possibility"?

Thanks!


\hfill

\texttt{Cologero on 2011-11-08 at 00:01 said: }

We mean the difference between essence and existence as explained here.


\hfill

\texttt{EXIT on 2011-11-10 at 10:24 said: }

Jesus/Satan: Will you lie, cheat and steal for me?

Dupe: I will.

Jesus/Satan: Will you play the race card?

Dupe: Yes, master.

Jesus/Satan: Will you follow my one and only messenger on earth, Cologero Salvo who knows better than all the perennial writers that ever lived?

Dupe: Anything you say, master. Christianity is the bestest, massa.


\hfill

\texttt{Charlotte Cowell on 2011-11-10 at 17:47 said: }

exit, I know this comment isn't directed at me but I'm just wondering how it is you personally came to hate Christianity / Jesus so much (or at least it seems that way). I'm not being sarcastic or trying to trip you up in some way, I'm genuinely wondering, there must be sadness behind it somewhere and I'm sorry (that isn't meant to be patronising, I'm not like that). Hope you are well. Cx


\hfill

\texttt{Charlotte Cowell on 2011-11-10 at 17:49 said: }

by the way I don't think you are wrong to point out that the enemy so often tries to lead us astray in the most subtle of ways, even by masquerading as the Master Jesus or the heavenly angels, it's a problem we all face in our lives (life) here on earth, tragically so for most people I would think


\hfill

\texttt{charlotte on 2011-11-12 at 15:58 said: }

i'd like to share this as Ysatis de Saint Simone has a very clear insight into true Christianity and also the era of the anti-Christ, her books are very interesting and help to understand the terrible confusions that have been caused because of the veiled language used in the Old Testament – those who understand the true and original Hebrew are few and far between, so it is no wonder so many are enraged and feel deceived. But we have to remember that rage and turning away from Christ is precisely what the bloodthirsty power of anti-Christ – the warmonger of the world – is hoping for. the war is waged within us and without us, we have to choose sides.

\url{http://video.google.com/videoplay?docid=-1291779146253225433}


\hfill

\texttt{Enda Miller on 2011-11-26 at 14:49 said: }

This post has had a dramatic impact on me, thank you very much! May I continue my questioning?

``For Guenon, Being is not the Infinite, which contains all possibilities. Being consists of all possibilities that are manifested. Non-being, then, consists of non-manifested possibilities as well as possibilities of non-manifestation."

Is this to suggest that for Guenon God = Infinite, that he understands the Infinite as the conflation of Being and Non-Being?


\hfill

\texttt{Cologero on 2011-11-27 at 00:53 said: }

Guenon understands the Infinite as having, in some sense, all possibilities. Things that exist (have being) are still possibilities, but with the property of manifestation, so from a logical point of view, you could say the Infinite is the union of Non-Being and Being.

Guenon did not explicitly claim God=Infinite. Rather, he said that theological language (God) needs to relate back to metaphysical language (Infinity). Hence, we filled in the details of that project by showing how we can understand God consistently, given Guenon's explanation of the Infinite. We drew on the best minds in the Western Tradition to show how everything that can be attributed to the Infinite has a correspondence to what we can say about God. Obviously, that is far from exhausting the topic, but it is a necessary first step in reclaiming the Western Tradition.


\hfill

\texttt{Ola on 2022-03-14 at 11:22 said: }

``From the strictly human point of view, which alone is what religions as such have in view, ``God" could not be the Absolute in itself, for the Absolute has no interlocutor; we may, however, say that God is the hypostatic Face turned towards the human world, or towards a particular human world; in other words, God is Divinity that personalizes itself in view of man and insofar as it takes on, to a greater or lesser extent, the countenance of a particular humanity." – F. Schuon, In the Face of the Absolute

``Although it is usually held that Plotinus inherited two major problems from his predecessors, i.e., a contradiction between the Pythagorean doctrine of the first Principle as an ultimate unity (the One) and the Peripatetic doctrine (going back to Anaxagoras) that the first Principle is the divine Intellect thinking itself, these two perspectives are in fact easily reconcilable, as numerous examples of different ancient theological systems prove.

In fact, only those who are bound to formal discursive reasoning can see a contradiction between Brahma nirguna and Brahma saguna, the ineffable Principle and personal Ishwara, Parama Shiva and Aham, Nun and Atum, or between apophatic and cataphatic ways of approaching God." – Algis Uzdavinys, The Heart of Plotinus

\end{sffamily}\end{footnotesize}
